\chapter{Testing and Debugging}
\label{ch:testing-debugging}

Code that is never tested is code you only \emph{hope} works. Tests turn hope into
evidence: small programs that check your code does what you claim, and keep
checking as you change it. And when something is wrong, debugging is how you find
out why. This chapter covers both writing tests with the \code{testing} module,
and the tools and habits that make finding bugs faster.

\section{Why test?}

A test is a tiny experiment: ``if I call this function with these inputs, I should
get this result.'' Writing tests has three payoffs. They catch mistakes
immediately, while the code is fresh in your mind. They guard against
\emph{regressions} a change that fixes one thing while quietly breaking
another because you can re-run every test after every change. And they document
how your code is meant to be used, in the most honest form possible: working
examples.

\section{Assertions with the \texttt{testing} module}

The \code{testing} module provides \textbf{assertions} statements that check a
condition and record a pass or a failure. You write a normal program that calls a
function and asserts the expected result:

\begin{salam}
import "testing"

func add(a : int, b : int) : int:
    ret a + b
end

func main:
    testing.AssertEqInt(add(2, 3), 5)
    testing.AssertEqInt(add(-1, 1), 0)
    testing.AssertTrue(add(2, 2) == 4)
    testing.Summary()
end
\end{salam}

\begin{progout}
tests: 3 passed, 0 failed
\end{progout}

Each \code{Assert} checks its condition and tallies the result; \code{testing.\-
Summary()} prints the totals at the end. When everything passes, you get a clean
report like the one above quiet confidence that your code does what you said.

\subsection{The assertion functions}

\begin{center}
\begin{tabular}{@{}ll@{}}
\toprule
\textbf{Assertion} & \textbf{Passes when} \\
\midrule
\code{testing.AssertTrue(c)}        & \code{c} is \code{true} \\
\code{testing.AssertFalse(c)}       & \code{c} is \code{false} \\
\code{testing.AssertEqInt(got, want)} & two integers are equal \\
\code{testing.AssertEqStr(got, want)} & two strings are equal \\
\code{testing.AssertEqBool(got, want)} & two booleans are equal \\
\code{testing.Assert(c)}            & \code{c} is \code{true} (general) \\
\bottomrule
\end{tabular}
\end{center}

The \code{...Eq...} forms are worth preferring over a plain \code{AssertTrue(a ==
b)} because, when they fail, they tell you \emph{both} values what you got and
what you wanted which is exactly what you need to diagnose the problem.

\subsection{Reading a failure}

A failing assertion does not stop the program; it records the failure, prints a
message, and lets the remaining checks run, so one run tells you about \emph{all}
the problems. Watch what happens when an expected value is wrong:

\begin{salam}
import "testing"

func add(a : int, b : int) : int:
    ret a + b
end

func main:
    testing.AssertEqInt(add(2, 3), 5)
    testing.AssertEqInt(add(2, 2), 5)    // wrong on purpose
    testing.Summary()
end
\end{salam}

\begin{progout}
FAIL: ints differ: got 4 want 5
tests: 1 passed, 1 failed
\end{progout}

The message names exactly what went wrong \code{got 4 want 5} and the
summary counts one pass and one failure. In a real test file you would fix the
code (or the expectation) until the summary reads all passes.

\section{Organising tests}

Group related checks into named functions, then call them from \code{main}. This
keeps each function focused on one piece of behaviour and makes the output easier
to follow:

\begin{salam}
import "testing"

func is_even(n : int) : bool:
    ret n % 2 == 0
end

func test_is_even():
    testing.AssertTrue(is_even(4))
    testing.AssertFalse(is_even(7))
    testing.AssertTrue(is_even(0))
end

func main:
    test_is_even()
    testing.Summary()
end
\end{salam}

\begin{progout}
tests: 3 passed, 0 failed
\end{progout}

As your program grows, keep test functions next to the code they exercise, give
them clear \code{test\_} names, and call them all from a single \code{main} that
ends in \code{testing.Summary()}. The exit code reflects the number of failures,
so a test program also slots neatly into an automated build that should stop when
a test fails.

\begin{tip}
Write a test the moment you find a bug, \emph{before} you fix it: a test that
reproduces the bug. Then fix the code until the test passes. Now you have not only
solved the problem but guaranteed it can never silently come back the test will
catch it if it does. This habit, fixing bugs ``test first,'' steadily builds a net
that makes your whole program safer to change.
\end{tip}

\section{Debugging: finding out why}

When a program misbehaves, the goal is to replace guessing with \emph{seeing}. The
most powerful debugging tool is also the humblest.

\subsection{Print debugging}

A well-placed \code{println} that shows a variable's value at a key moment will
solve the large majority of bugs. When a result is wrong, print the inputs and
intermediate values along the path that produces it, and compare what you see with
what you expected. The discrepancy is almost always where the bug lives:

\begin{salam}
func average(a : int[3]) : f64:
    mut sum : int = 0
    each x in a:
        sum += x
        println "after adding", x, "sum is", sum    // a temporary probe
    end
    ret (sum as f64) / 3.0
end

func main:
    data : int[3] = [10, 20, 30]
    println "average:", average(data)
end
\end{salam}

\begin{progout}
after adding 10 sum is 10
after adding 20 sum is 30
after adding 30 sum is 60
average: 20
\end{progout}

The probe lines confirm the running total is building correctly. Once you have
found and fixed the problem, remove the temporary prints they were scaffolding,
not part of the finished program.

\subsection{Let the compiler help}

Many bugs never reach run time at all: Salam's static checks catch them first. A
type mismatch, an immutable variable you tried to reassign, an unused parameter, an
undefined name each is reported at compile time with a clear message and the
line it occurred on. Read those messages carefully; they are the cheapest bug
reports you will ever get, arriving before the program even runs. Much of what
would be a debugging session in a dynamic language is, in Salam, simply a compiler
error you fix on the spot.

\subsection{Heavier tools}

For the stubborn cases, the compiler offers more. Building with debug information
(\code{salam debug}) lets you step through your program in a debugger, inspecting
variables as it runs. And because Salam compiles to C, the memory checker
(\code{salam memcheck}, which builds with AddressSanitizer) pinpoints
use-after-free and out-of-bounds errors invaluable when you have been working
with manual memory or the FFI. You will not need these often, but they are there
when a print statement is not enough.

\begin{deep}[In Depth: a debugging method]
When a bug resists, work like a scientist rather than thrashing. \textbf{Reproduce}
it reliably find the smallest input that triggers it. \textbf{Localise} it by
bisection: confirm the data is right at the start and wrong at the end, then check
the middle, narrowing the region until you find the line where right becomes wrong.
\textbf{Form a hypothesis} about the cause, and \textbf{test} it with a targeted
print or a change. The smallest reproducing case plus bisection will corner almost
any bug, far faster than staring at the code hoping for inspiration.
\end{deep}

\section{Chapter summary}

\begin{itemize}[leftmargin=1.4em]
  \item Tests turn ``I hope it works'' into evidence and guard against
        regressions.
  \item The \code{testing} module asserts conditions (\code{AssertTrue},
        \code{AssertEqInt}, \code{AssertEqStr}, \dots) and reports totals with
        \code{Summary}; failures print both values and do not halt the run.
  \item Group checks into \code{test\_}-named functions called from \code{main}.
  \item Write a failing test that reproduces a bug before fixing it.
  \item \code{println} probes solve most bugs; remove them once done.
  \item Salam's compile-time checks catch many bugs before run time read the
        messages.
  \item \code{salam debug} and \code{salam memcheck} handle the hard cases.
\end{itemize}

\section{Exercises}

\exercise Write a function \code{factorial(n int) int} and a \code{test\_factorial}
that asserts \code{factorial(0) == 1}, \code{factorial(5) == 120}, and one more
case of your choosing.

\exercise Add a deliberately wrong assertion to a passing test file and study the
failure message and the summary line.

\exercise Take a buggy function (for example, a loop with an off-by-one error) and
use \code{println} probes to locate the exact line where the values go wrong.

\exercise Write tests for a \code{clamp(x, lo, hi)} function covering a value below
the range, inside it, and above it.

\exercise Describe, in a comment, the difference between a bug Salam catches at
compile time and one you would only find by running tests, with an example of
each.
